% ============================================================
\section{Conclusion}

This thesis presented the design, implementation, and evaluation of a mixed-reality pipeline for six-degree-of-freedom (6D) object pose estimation with Apple Vision Pro (AVP) as the interactive front-end. The final system couples a native visionOS client with a Python backend that ingests RGB via AirPlay/UxPlay, acquires metric depth via an Intel RealSense sensor, reprojects depth into the UxPlay camera model using ArUco-based camera-to-camera registration, synthesizes ROI masks from normalized parameters, and orchestrates inference through a Dockerized FoundationPose REST service. This split architecture isolates user interaction and rendering on the headset from pixel-level processing and GPU inference on external hardware.

A central constraint throughout development was operation under the standard Apple Developer Program (ADP). Without ADEP entitlements, direct access to AVP camera frames, depth sensing, and native intrinsics is unavailable. This constraint invalidated an initial AVP-sensor-first design and forced a deployable alternative based on mirrored RGB acquisition and external RGB-D sensing. As a consequence, calibration, registration, mask synthesis, and inference orchestration were implemented on the backend, while the visionOS application remained intentionally lightweight, providing ROI interaction, model selection, device/head pose streaming, and pose visualization.

Implementation followed a staged workflow. Early iterations used a Python prototype to validate calibration, ROI handling, depth sources, and request payload construction before committing to the final visionOS UI. During integration, a Streamlit-based testing frontend and a deterministic request replay script were used to decouple pose inference validation from end-to-end headset integration and to reproduce failures reliably. The final visionOS client was implemented with SwiftUI, ARKit, and RealityKit, supporting windowed configuration and immersive overlays, with pose outputs rendered as lightweight primitives to keep the client responsive.

Evaluation was structured to separate camera-frame pose consistency from world-space alignment error. Two paired pipelines were compared: a RealSense reference pipeline that estimates pose directly in the RealSense color camera frame, and an RS$\rightarrow$UxPlay pipeline that estimates pose in the UxPlay camera frame using reprojected RealSense depth, then transforms results back into the RealSense frame for comparison. This design isolates error contributions introduced by registration and reprojection when using the mirrored view, while treating world-space overlays as an outlook metric dominated by additional factors such as ARKit tracking drift, intrinsics/distortion uncertainty, ArUco estimation noise, and sensitivity to the device-to-camera offset. The results confirm that interactive pose feedback is achievable under current platform constraints, while highlighting the dominant practical limitations: reprojection artifacts in transformed depth, drift in immersive overlays over time, and sensitivity of world-space alignment to calibration quality. Temporal robustness strategies for occlusion and outlier handling were investigated, but were not enabled by default during evaluation to avoid masking upstream alignment issues.

Overall, the thesis demonstrates that AVP can act as an effective mixed-reality control and visualization layer for external, AI-based 6D pose estimation even under restricted native sensor access. The system provides a concrete basis for extending mixed-reality interaction toward AR-assisted Programming by Demonstration workflows, while making the remaining bottlenecks—registration robustness, calibration stability, and drift handling—explicit and measurable.

\subsection{Outlook}

Future improvements follow directly from the identified constraints and failure modes. First, any future availability of native camera/depth access or equivalent sensing interfaces would remove the need for AirPlay capture and external depth, reducing calibration overhead and eliminating an entire reprojection stage. Second, stability can be improved by strengthening registration: increasing ArUco target quality and observability, adding quality gating and temporal averaging for $\mathbf{T}_{AVP \leftarrow RS}$, and replacing manual device-to-camera offset tuning with a guided, quantitative calibration procedure based on reprojection error. Third, transformed depth artifacts can be mitigated by ROI-aware validity masks and optional hole handling so that invalid reprojection regions are not forwarded to inference. Fourth, world-space drift can be reduced by periodic re-anchoring to a physical reference and by surfacing anchor-quality cues in the UI to trigger re-calibration when needed. Finally, the interaction workflow can be hardened for real usage by adopting snapshot-style ROI capture to avoid recursive mirroring contamination and by validating the multi-window interface in structured usability studies.

% These directions preserve the core architectural idea of the thesis—thin mixed-reality interaction on AVP combined with external perception and inference—while targeting the dominant sources of residual error observed during evaluation.
